\documentclass{article}

\usepackage{amsmath,amssymb}
\usepackage{xurl}
\usepackage[hidelinks]{hyperref}
\setlength{\emergencystretch}{2em}

% Shared commands for notes in docs/paper-gaps/.

\DeclareUnicodeCharacter{2080}{\ifmmode{}_{0}\else\textsubscript{0}\fi}
\DeclareUnicodeCharacter{2081}{\ifmmode{}_{1}\else\textsubscript{1}\fi}
\DeclareUnicodeCharacter{2082}{\ifmmode{}_{2}\else\textsubscript{2}\fi}
\DeclareUnicodeCharacter{2099}{\ifmmode{}_{n}\else\textsubscript{n}\fi}

\newcommand{\Lean}{\textsc{Lean}}
\ExplSyntaxOn
\NewDocumentCommand{\leanid}{m}{
  \ifmmode
    \textsf{\detokenize{#1}}
  \else
    \group_begin:
    \urlstyle{sf}
    \tl_set:Nn \l_tmpa_tl {#1}
    \tl_replace_all:Nnn \l_tmpa_tl {\_} {_}
    \exp_args:NV \path \l_tmpa_tl
    \group_end:
  \fi
}
\ExplSyntaxOff
\providecommand{\tr}{\operatorname{tr}}
\newcommand{\tnleanIssueBase}{https://github.com/LionSR/TNLean/issues/}
\newcommand{\tnleanPrBase}{https://github.com/LionSR/TNLean/pull/}
\newcommand{\tnleanDocBase}{https://sirui-lu.com/TNLean/docs/}
\newcommand{\ghissue}[1]{\href{\tnleanIssueBase#1}{GitHub issue \##1}}
\newcommand{\ghpr}[1]{\href{\tnleanPrBase#1}{PR \##1}}
\newcommand{\doclink}[1]{\href{\tnleanDocBase#1.html}{\path{#1}}}

% Dirac notation and the identity operator, as in blueprint/src/macros/common.tex.
% \providecommand keeps note-local definitions authoritative.
\usepackage{dsfont}
\providecommand{\ket}[1]{|#1\rangle}
\providecommand{\bra}[1]{\langle #1|}
\providecommand{\braket}[2]{\langle #1 | #2 \rangle}
\providecommand{\ketbra}[2]{|#1\rangle\!\langle #2|}
\providecommand{\Id}{\mathds{1}}
\providecommand{\one}{\mathds{1}}
\providecommand{\C}{\mathbb{C}}
\providecommand{\MN}[1]{M_{#1}(\C)}
\providecommand{\MD}{M_D(\C)}

% Verdict marker: \gapnote{<kind>}{<status>}. Kind is the policy
% classification (clarification, local-correction, scope-restriction,
% unfaithful, false-source, open-gap); status is open, wip (elimination
% underway), resolved, or historical. Machine-read by the site index;
% prints a verdict line.
\newcommand{\gapnote}[2]{\par\noindent\textbf{Verdict:} #1 (#2).\par}


% Fallbacks keep this note compilable when a different command.tex is found first.
\providecommand{\leanid}[1]{\textsf{\detokenize{#1}}}
\providecommand{\gapnote}[2]{\par\noindent\textbf{Verdict:} #1 (#2).\par}

\title{Two Factors of Two in the Doubling Iteration of Hastings' Area Law}
\author{}
\date{2026-10-28}

\begin{document}
\maketitle
\gapnote{local-correction}{resolved}

\section{The source passage}

All entropies in this note are the conventional von Neumann entropy
$S(\rho)=-\operatorname{tr}(\rho\ln\rho)$. The definition immediately
before the main theorem in the source omits the minus sign. The subsequent
use of $0\le S_l\le l\ln D$ and of entropy subadditivity fixes the intended
convention; the formalization uses this conventional nonnegative entropy.
This typographical correction adds no hypothesis.


In the proof of the one-dimensional area law~\cite{Hastings2007AreaLaw}
(arXiv:0705.2024v4, lines~197--231), $S_l$ is the largest entropy of an
interval of $l$ sites inside a window of $l_0$ sites. The proof uses
$0\le S_l\le l\ln D$ and the strengthened doubling inequality
\begin{align}
    S_{2l}\le 2S_l-\bigl(1-2C_1e^{-l/\xi'}\bigr)\frac{l}{\xi'}+\ln C_1+C_2 ,
    \label{eq:hastings07_doubling_constants_claim}
\end{align}
sets $\xi_0=2\xi'\ln(2C_1)$, and iterates from $l=\xi_0$ to
$l=2\xi_0,4\xi_0,\dots$. Lines~221--224 assert that for $l>\xi_0$
\begin{align}
    S_l\le l\ln D-\frac{l\lfloor\log_2(l/\xi_0)\rfloor}{\xi'}
    +\bigl(2+\ln C_1+C_2\bigr)\frac{l}{\xi_0},
    \label{eq:hastings07_doubling_constants_source_bound}
\end{align}
and lines~225--231 conclude
\begin{align}
    l_0\le\xi'\ln(2C_1)\,
    2^{\lceil\xi'\ln D+(2+C_2)\xi'/\xi_0+1/2\rceil}.
    \label{eq:hastings07_doubling_constants_source_window}
\end{align}

\section{The two corrections}

Dividing~\eqref{eq:hastings07_doubling_constants_claim} by $2l$ shows that one
doubling lowers $S_l/l$ by
$1/(2\xi')$ up to the error terms, not by $1/\xi'$. After $N$ doublings
from a length $l_1\ge\xi_0$ the correct bound at $l=2^Nl_1$ is
\begin{align}
    S_l\le l\ln D-\frac{lN}{2\xi'}+\bigl(2+\ln C_1+C_2\bigr)\frac{l}{\xi_0}.
    \label{eq:hastings07_doubling_constants_corrected_bound}
\end{align}
The error constant $2$ survives: the exponential errors contribute at most
$l/\xi_0$, since $\sum_{n\ge0}e^{-2^nt}\le e^{-t}/(1-e^{-t})$ for
$t=l_1/\xi'\ge2\ln(2C_1)$ and $x\ln x\le x^2-1$ for $x=2C_1\ge1$, and the
constant terms contribute at most $(\ln C_1+C_2)\,l/\xi_0$ by the geometric
sum $\sum_n2^{-n}\le2$. The function
$f(l)=l\bigl(\ln D-\log_2(l/l_1)/(2\xi')\bigr)$ satisfies
$f(2l)=2f(l)-l/\xi'$ exactly, so the printed decrease
$l\lfloor\log_2(l/\xi_0)\rfloor/\xi'$ cannot follow
from~\eqref{eq:hastings07_doubling_constants_claim}; the same factor is missing
from the heuristic
$S_l\approx l\ln D-\log_2(l/\xi_0)\,l/\xi'$ at lines~212--213.
On the lengths $l=2^nl_1$ that the iteration visits, $f$ also satisfies
$f(l)\le l\ln D$, and $f(l)\ge0$ for $l\le l_1 4^{\xi'\ln D}$; there it
meets the claim with every error term dropped. The iteration, which uses the
claim only at these lengths, therefore yields $4^{\xi'\ln D}$ and no better;
the printed $2^{\xi'\ln D}$ would need the decrease $2l/\xi'$
in~\eqref{eq:hastings07_doubling_constants_claim}, or an argument that uses the
claim at other lengths.

The bound becomes negative at $N=M$ with
\begin{align}
    M=\max\Bigl\{0,\Bigl\lceil2\Bigl(\xi'\ln D+\frac{(2+C_2)\xi'}{\xi_0}
    +\frac12\Bigr)\Bigr\rceil\Bigr\},
    \label{eq:hastings07_doubling_constants_exponent}
\end{align}
twice the printed exponent, because $\ln C_1\,\xi'/\xi_0<1/2$. The
contradiction at length $l_12^M$ gives $l_0<l_12^M$. The printed prefactor
$\xi'\ln(2C_1)=\xi_0/2$ is half the starting length; the correct prefactor
is the starting length itself.

Lengths are natural numbers, so the formalization starts at
$l_1=\lceil\xi_0\rceil$ and proves
\begin{align}
    l_0<\lceil\xi_0\rceil\,2^M .
    \label{eq:hastings07_doubling_constants_corrected_window}
\end{align}
In the main theorem these corrections replace $2^{\xi'\ln D}$ in $S_{\max}$
by $4^{\xi'\ln D}$; the form of the area law is unchanged.

\section{Implicit hypotheses}

The source uses three facts left implicit, all recorded as hypotheses.
\begin{itemize}
    \item $2C_1>1$, so that $\xi_0>0$ is a length.
    \item $\ln C_1+C_2\ge0$, used when the constant errors
          $(\ln C_1+C_2)/(2\cdot2^nl_1)$ are bounded by
          $(\ln C_1+C_2)/\xi_0$. Increasing $C_2$
          keeps~\eqref{eq:hastings07_doubling_constants_claim} true, so this
          costs nothing.
    \item Inequality~\eqref{eq:hastings07_doubling_constants_claim} is stated
          ``for all $l\le l_0$''; it is used only when both intervals fit in
          the window, $2l\le l_0$, and is assumed only there.
\end{itemize}
No hypothesis on $D$ is needed.

\section{Choosing another Lieb--Robinson length}

The correction concerns the iteration at a fixed decay length. It does not
prevent an entropy estimate of the printed form after another choice of
the velocity and length in the Lieb--Robinson bound. Those constants are
chosen existentially in the source, lines~96--112.

Let $\zeta$ be the length of the proved projector factorization, and let
$C_1=\max\{K\zeta^3,1\}$ with $K$ depending only on the interaction range.
Set $x=6\zeta$ and $a=\ln(K+1)+3$. The corrected iteration, with its
bootstrap remainder, gives
\begin{align}
    S([0,i))&\le B(a)x\ln x\ln D\,2^{2x\ln D},
    \label{eq:hastings07_doubling_constants_polynomial_bound}
    \\
    B(a)&=6(2a+5)2^{9/\ln2+1}+2a+20.
    \label{eq:hastings07_doubling_constants_coefficient}
\end{align}
For every $s\ge2\zeta$, put $\xi'=6s$. Since $\xi'\ge2x$, increasing the
positive factors in
(\ref{eq:hastings07_doubling_constants_polynomial_bound}) gives
\begin{align}
    S([0,i))\le B(a)\xi'\ln\xi'\ln D\,2^{\xi'\ln D}.
    \label{eq:hastings07_doubling_constants_enlarged_bound}
\end{align}
The constant depends only on the interaction range and is numerical in the
nearest-neighbour case. The case $D=1$ is immediate because its entropy
vanishes.

There is a choice of $s$ of the source's form with velocity proportional
to $J$. With $c_R=2(R+1)2^{2R+1}e^{2R}$, take
\begin{align}
    v=40c_RJ,\qquad \xi_C=80,\qquad
    s=\max\{2v/\Delta E,\xi_C\}.
    \label{eq:hastings07_doubling_constants_source_pair}
\end{align}
For $J>0$ this is an admissible pair in the reciprocal-length convention
of \path{hastings07_lemma1_decay_length.tex}: the commutator has numerical
prefactor $4$, arbitrary finite supports, and decay $e^{-d/\xi_C}$ on
$|t|\le d/v$. Since $\zeta=20(1+c_RJ/\Delta E)$, the elementary inequality
$2(1+u)\le4\max\{u,1\}$ for $u\ge0$ yields $2\zeta\le s$.
Thus (\ref{eq:hastings07_doubling_constants_enlarged_bound}) has the
printed dependence on the new length. It does not validate the smaller
fixed-length decrement or window estimate in the source.

\section{Formal statements}

The bound on $S_l$ is
\leanid{TNLean.MPS.AreaLaw.Hastings.le_of_doubling} and the bound on $l_0$
is \leanid{TNLean.MPS.AreaLaw.Hastings.lt_of_doubling}.
The entropy estimate
(\ref{eq:hastings07_doubling_constants_enlarged_bound}) is
\leanid{TNLean.MPS.AreaLaw.Hastings.exists_regionEntropy_prefix_le_printed}.
The comparison $2\zeta\le s$ is
\leanid{TNLean.MPS.AreaLaw.two_mul_lemmaOneEndpointLength_le_sourceLength}.

\bibliographystyle{alpha}
\bibliography{references}

\end{document}
